% ===========================================================================
% Classic-roofline version of Section 3 (replaces roofline.tex, which is kept on disk).
% Drop-in: reuses every label referenced from sections 04/05.
% Assumptions are typeset in red via \assume{...}.
% ===========================================================================
\providecommand{\assume}[1]{{\color{red}#1}}

\section{A Communication-Centric Roofline}
\label{sec:roofline}
\label{sec:simulator}  % keeps existing \ref{sec:simulator} links alive

\paraq{What is this section for?}
Section~\ref{sec:hybrids} showed that the fastest strategy---data parallelism, gsplit, or a
hybrid---changes with the model, the graph, the depth, the hidden size, and the interconnect.
This section places every layer of every strategy on a classic roofline. The workload (model,
graph, partition, strategy) fixes where a point sits, and the hardware (GPU and interconnect)
fixes where the roofs are. Changing the GPU therefore redraws the roofs over the same points, and
the bottleneck of each layer can be read off for any machine. Section~\ref{sec:predicted} plots the
roofline for representative settings, and Section~\ref{sec:results} compares its picture with the
measurements.

\paraq{The model is deliberately simple: a first iteration that guides data collection.}
The roofline is built from the fewest quantities that still separate the strategies: two numbers
per layer, both obtained from sampling, and three measured rates for each GPU, interconnect, and
hidden size. The plan has three steps. First, place every layer of every strategy on the roofline
for a set of representative settings. Second, collect those points together with the measured
winner of each setting. Third, inspect the plots for the patterns that separate the winners, and
turn them into an algorithm that picks the hybrid strategy from the roofline alone. Assumptions
that the measurements show to matter are refined in later iterations; until then they are listed
in red.

% ---------------------------------------------------------------------------
\subsection{Roofline assumptions}
\label{sec:cm-purpose}

\paraq{Which classic roofline assumptions does the model adopt?}
The roofline asks which resource limits a layer, not how long the layer takes. Three resources
are modelled: dense arithmetic, scattered DRAM reads (the neighbour gather), and the interconnect.
The model adopts the standard roofline assumptions:
\begin{enumerate}\itemsep1pt
  \item[\assume{(R1)}] \assume{\textbf{Perfect overlap.} All resources run concurrently, so the
        time of a layer is the maximum of the times each resource needs, and the attainable
        throughput is the minimum over the roofs.}
  \item[\assume{(R2)}] \assume{\textbf{Steady state.} Latency, kernel-launch overhead,
        synchronisation, and pipeline fill and drain are ignored.}
  \item[\assume{(R3)}] \assume{\textbf{Fixed ceilings.} Each resource has one sustained rate per
        hidden size, measured once by a microbenchmark at the configured fanout and independent of
        the layer.}
  \item[\assume{(R4)}] \assume{\textbf{Compulsory traffic.} Bytes are counted from the algorithm:
        each pulled row crosses the interconnect once and each neighbour row is gathered once per
        edge; caches are ignored.}
  \item[\assume{(R5)}] \assume{\textbf{One kernel per point.} Each layer is analysed in isolation;
        successive layers add, and the backward pass scales every resource by the same factor and
        is dropped.}
\end{enumerate}
Under (R1)--(R5) the roofline is an upper bound on throughput. It is used to identify the binding
resource of each layer, not to predict time.

% ---------------------------------------------------------------------------
\subsection{Work of one node}
\label{sec:cm-node}

\paraq{What work does computing a single node involve?}
The hidden representation $H^{l}_u$ of a destination node $u$ in layer $l$ is computed from the
hidden representations of its sampled neighbours, $H^{l-1}_v$ for $v \in N(u)$, and the layer's
weights. (In these equations $l$ runs in the forward direction; from
Section~\ref{sec:cm-axes} on, layers are indexed by $\ell$, counted from the seeds.) For
GraphSAGE (mean aggregator):
\begin{align}
  H^{l}_u &= \sigma\Big( W^{\mathrm{self}}_l H^{l-1}_u
      + W^{\mathrm{neigh}}_l \tfrac{1}{|N(u)|}\textstyle\sum_{v \in N(u)} H^{l-1}_v \Big).
      \label{eq:sage}
\end{align}
For GAT:
\begin{align}
   z^{l}_v &= W_l\, H^{l-1}_v,
      \qquad v \in N(u) \cup \{u\}
      && \text{dense: one GEMM row per source node} \label{eq:gat-dense}\\
   \alpha_{uv} &= \underset{v \in N(u)}{\mathrm{softmax}}\Big(
         \mathrm{LeakyReLU}\big(a_l^{\top}\,[\, z^{l}_u \,\|\, z^{l}_v \,]\big)\Big)
      && \text{sparse: one scalar score per edge} \label{eq:gat-attn}\\
   H^{l}_u &= \sigma\Big( \sum_{v \in N(u)} \alpha_{uv}\, z^{l}_v \Big)
      && \text{sparse: weighted gather over } N(u) \label{eq:gat-agg}
\end{align}
Both models interleave \emph{dense} work (vector--matrix and vector--vector products over
contiguous memory, bound by the arithmetic rate) with \emph{sparse} work (the aggregation over
$N(u)$, which reads neighbour rows at scattered addresses and is bound by gather bandwidth).

\paraq{How are the dense and sparse work of a node counted?}
Every hidden layer has width $h$; a hidden row is $s_H = h\,b$ bytes with $b$ bytes per value.
For a node with $e$ sampled in-edges, a model is described by five constants read off its layer
equations (Table~\ref{tab:cm-constants}):
\begin{itemize}\itemsep1pt
  \item $\kappa$: hidden rows gathered per edge;
  \item $\beta_e, \beta_n$: vector--matrix products ($h^2$ multiply--adds each) per edge and per
        node;
  \item $\gamma_e, \gamma_n$: vector--vector products ($h$ multiply--adds each) per edge and per
        node.
\end{itemize}
Counting a multiply--add as two FLOPs,
\begin{align}
  F(e) &= 2\big[(\beta_e\, e + \beta_n)\, h^{2} + (\gamma_e\, e + \gamma_n)\, h\big]
       \quad\text{FLOPs}, \label{eq:dense-cost}\\
  S(e) &= \kappa\, e\, s_H \quad\text{bytes gathered.} \label{eq:sparse-cost}
\end{align}
\assume{All hidden layers have width $h$; input-feature loading and the first layer's feature
width are not modelled. The per-edge dense products are charged to every destination, which
over-counts GAT, whose $z_v$ is computed once per source row.}

\begin{table}[h]
\centering
\caption{Model constants of Eqs.~\eqref{eq:dense-cost}--\eqref{eq:sparse-cost} (forward pass,
per destination node). \placeholder{} = to be derived from the layer equations.}
\label{tab:cm-constants}
\small
\begin{tabular}{l c c c c c}
\toprule
Model & $\kappa$ & $\beta_e$ & $\beta_n$ & $\gamma_e$ & $\gamma_n$ \\
\midrule
SAGE & \placeholder & \placeholder & \placeholder & \placeholder & \placeholder \\
GAT  & \placeholder & \placeholder & \placeholder & \placeholder & \placeholder \\
\bottomrule
\end{tabular}
\end{table}

% ---------------------------------------------------------------------------
\subsection{The axes}
\label{sec:cm-axes}

\paraq{Which quantities does sampling provide for one layer?}
A strategy fixes which rows each GPU computes and which rows it receives, and one sampling pass
reveals both, without features or training. For strategy $s$, layer $\ell$ (counted from the
seeds, as in Section~\ref{sec:hybrid-dp}), and GPU $q$, sampling gives
\begin{itemize}\itemsep1pt
  \item $n_{\ell,q}$: destination rows computed;
  \item $\bar{e}_{\ell,q}$: their average fanout;
  \item $m_{\ell,q}$: rows received over the interconnect ($m=0$ for data parallelism);
  \item $u_\ell$: distinct destination rows of the layer over all $G$ GPUs.
\end{itemize}
The shuffle synchronises the GPUs, so a layer ends when its slowest GPU does. Each layer is
therefore represented by its \textbf{busiest GPU} $q^{*}$, the GPU with the longest time under (R1)
(Section~\ref{sec:cm-point}). \assume{Representing the layer by its busiest GPU ignores how the
other GPUs are loaded.} Below, $n$, $\bar{e}$, $m$ denote the counts of $q^{*}$ in layer $\ell$.

\paraq{The $x$-axis is the operational intensity with respect to the interconnect.}
The $x$-coordinate of a layer is the number of FLOPs its busiest GPU performs per byte it receives
over the interconnect,
\begin{equation}
  I_\ell \;=\; \frac{n\, F(\bar{e})}{m\, s_H} \quad \text{FLOP/byte}.
  \label{eq:cm-x}
\end{equation}
This is the classic arithmetic intensity, with the interconnect in place of DRAM. It depends only
on the model, the hidden size, and the sampled blocks of the strategy, not on the GPU. A
data-parallel layer receives nothing, so its intensity is infinite.

\paraq{The $y$-axis is the useful FLOP rate per GPU.}
Data parallelism computes some rows on several GPUs, and every strategy can be imbalanced. Both
effects are captured by the \emph{work inflation}
\begin{equation}
  \delta_\ell \;=\; \frac{G\, n}{u_\ell} \;\ge\; 1,
  \label{eq:cm-redundancy}
\end{equation}
the ratio of the rows the busiest GPU computes to its fair share $u_\ell/G$. The $y$-coordinate
counts only useful FLOPs, those of the fair share, so a layer that attains $A$ FLOP/s plots at
$y = A/\delta_\ell$. Like $I_\ell$, $\delta_\ell$ comes from sampling alone.

% ---------------------------------------------------------------------------
\subsection{The roofs}
\label{sec:cm-roofs}

\paraq{The hardware contributes three roofs.}
Let $\hat{P}$ be the measured dense rate, $\hat{B}_{\mathrm{mem}}$ the measured gather bandwidth,
and $\hat{B}_{\mathrm{net}}$ the measured all-to-all bandwidth (Table~\ref{tab:cm-microbench}).
For a layer of intensity $I$, the attainable FLOP rate is bounded by
\begin{itemize}\itemsep1pt
  \item the \textbf{dense roof} $\hat{P}$, flat;
  \item the \textbf{gather roof} $\hat{B}_{\mathrm{mem}}/g$, flat, where
        $g = S(\bar{e})/F(\bar{e})$ is the bytes gathered per FLOP;
  \item the \textbf{interconnect roof} $\hat{B}_{\mathrm{net}}\, I$, a line of slope $+1$ on
        log--log axes.
\end{itemize}
The two flat roofs together form the \emph{compute ceiling}
\begin{equation}
  C \;=\; \min\!\Big(\hat{P},\; \frac{\hat{B}_{\mathrm{mem}}}{g}\Big).
  \label{eq:roof-cmp}
\end{equation}
The gather roof depends on the model and hidden size through $g$; this is how GAT and SAGE obtain
different attainable FLOP rates on the same GPU. When the $h^2$ term of $F$ dominates,
$g \propto 1/h$: small hidden sizes tend to be gather-bound and large ones dense-bound.

\paraq{The attainable throughput is the minimum of the roofs.}
By (R1), a layer runs at the lowest roof that applies to it:
\begin{equation}
  y \;=\; \frac{1}{\delta}\,\min\!\big( C,\; \hat{B}_{\mathrm{net}}\, I \big).
  \label{eq:cm-roofline}
\end{equation}
The roofs meet at the \textbf{ridge}
\begin{equation}
  I^{*} \;=\; \frac{C}{\hat{B}_{\mathrm{net}}}.
  \label{eq:ridge}
\end{equation}
A layer with $I < I^{*}$ is \textbf{communication-bound}; a layer with $I > I^{*}$ is
\textbf{compute-bound}. A slower interconnect moves the ridge to the right, and so does a faster
GPU, which raises $C$.

\begin{table}[h]
\centering
\caption{Measured hardware rates (microbenchmarks on one A100-SXM4-80GB node, fp32, TF32 off;
$10^5$ destination rows, fanout $e{=}15$; all-to-all over 4 GPUs, bytes received per GPU).
These are the only hardware inputs of the roofline. Raw data: \texttt{results/microbench/}.}
\label{tab:cm-microbench}
\small
\begin{tabular}{l l c c c c c}
\toprule
Rate & Microbenchmark & $h{=}256$ & $512$ & $1024$ & $2048$ & $4096$ \\
\midrule
$\hat{P}$ (TFLOP/s) & one layer's GEMM on a sampled block
  & 15.0 & 18.7 & 19.0 & 19.1 & 19.1 \\
$\hat{B}_{\mathrm{mem}}$ (GB/s) & neighbour gather, fanout $e$ (CSR SpMM)
  & 1206 & 1240 & 1259 & 1263 & 1281 \\
\quad alt.\ kernel & index\_select $+$ sum
  & 225 & 254 & 197 & 196 & 184 \\
$\hat{B}_{\mathrm{net}}$ (GB/s), PCIe-emu & NCCL all-to-all of $s_H$-byte rows
  & 12.2 & 13.1 & 13.3 & 13.3 & 13.3 \\
$\hat{B}_{\mathrm{net}}$ (GB/s), NVLink & NCCL all-to-all of $s_H$-byte rows
  & 154 & 177 & 185 & 196 & 200 \\
\bottomrule
\end{tabular}
\end{table}

% ---------------------------------------------------------------------------
\subsection{Where each strategy lands}
\label{sec:cm-point}
\label{sec:cm-strategies}

\paraq{Each layer of each strategy is one point, computed from two times.}
On GPU $q$, layer $\ell$ needs compute time and transfer time
\begin{equation}
  T^{\mathrm{cmp}}_{\ell,q} = \frac{n_{\ell,q}\, F(\bar{e}_{\ell,q})}{C},
  \qquad
  T^{\mathrm{net}}_{\ell,q} = \frac{m_{\ell,q}\, s_H}{\hat{B}_{\mathrm{net}}}.
  \label{eq:cm-times}
\end{equation}
The busiest GPU is $q^{*} = \arg\max_q \max(T^{\mathrm{cmp}}_{\ell,q}, T^{\mathrm{net}}_{\ell,q})$,
and the layer's point is
\begin{equation}
  \big(I^{s}_\ell,\; y^{s}_\ell\big)
  \quad\text{with}\quad
  y^{s}_\ell = \frac{1}{\delta^{s}_\ell}\,\min\!\big( C,\; \hat{B}_{\mathrm{net}}\, I^{s}_\ell \big),
  \label{eq:cm-point}
\end{equation}
where $I^{s}_\ell$ and $\delta^{s}_\ell$ are Eqs.~\eqref{eq:cm-x} and~\eqref{eq:cm-redundancy}
evaluated on $q^{*}$. An $L$-layer model gives $L$ points per strategy.

\paraq{A split layer whose compute and shuffle overlap sits on the roof.}
A split layer computes each row once, so $\delta^{\mathrm{split}}_\ell$ reflects imbalance only and
is close to $1$. When compute and shuffle overlap, as (R1) assumes and as the pipeline hybrid
arranges across two micro-batches, the point lies on the roof:
\begin{equation}
  y^{\mathrm{pipe}} \;=\; \frac{1}{\delta^{\mathrm{split}}}\,
      \min\!\big( C,\; \hat{B}_{\mathrm{net}}\, I \big).
  \label{eq:cm-pipe-point}
\end{equation}
\assume{The pipeline is taken to realise (R1) exactly; micro-batch fill and drain and the lower
efficiency of half-size GEMMs are ignored.}

\paraq{gsplit runs compute and shuffle one after the other, so it sits below the roof.}
gsplit does not overlap the two, and its time is their sum:
\begin{equation}
  \frac{1}{y^{\mathrm{gsplit}}} \;=\; \delta^{\mathrm{split}}
      \Big(\frac{1}{C} + \frac{1}{\hat{B}_{\mathrm{net}}\, I}\Big).
  \label{eq:cm-gsplit-point}
\end{equation}
It is drawn directly below the pipeline point of the same layer. The vertical gap is
\begin{equation}
  \frac{y^{\mathrm{pipe}}}{y^{\mathrm{gsplit}}}
  \;=\; \frac{T^{\mathrm{cmp}} + T^{\mathrm{net}}}{\max(T^{\mathrm{cmp}}, T^{\mathrm{net}})}
  \;\in\; [1, 2],
  \label{eq:cm-pipe-gain}
\end{equation}
largest at the ridge, where $T^{\mathrm{cmp}} = T^{\mathrm{net}}$. Far from the ridge one
resource dominates and overlap has little to hide.

\paraq{A data-parallel layer is a flat line below the compute ceiling.}
A data-parallel layer receives nothing, so $I = \infty$ and the interconnect roof never binds. It
attains the compute ceiling but spends a fraction $1 - 1/\delta^{\mathrm{dp}}$ of it on redundant
rows:
\begin{equation}
  y^{\mathrm{dp}} \;=\; \frac{C}{\delta^{\mathrm{dp}}}.
  \label{eq:cm-dp-point}
\end{equation}
It is drawn as a horizontal line at that height. The computation trees of different GPUs overlap
more far from the seeds, so $\delta^{\mathrm{dp}}_\ell$ grows with $\ell$ and outer layers have
lower lines.

\paraq{How is the winner of a layer read off the plot?}
For layer $\ell$, compare the heights of its pipeline point, its gsplit point, and its
data-parallel line; the highest wins. Only heights of the same layer are compared, and layers are
distinguished by colour. Writing $\Delta_\ell = \delta^{\mathrm{dp}}_\ell /
\delta^{\mathrm{split}}_\ell$ for the redundancy ratio, Eqs.~\eqref{eq:cm-pipe-point},
\eqref{eq:cm-gsplit-point}, and~\eqref{eq:cm-dp-point} give
\begin{equation}
  \boxed{\;
  \begin{aligned}
    \text{data parallelism beats gsplit at layer } \ell
      &\iff \Delta_\ell \;<\; 1 + \frac{I^{*}}{I_\ell},\\
    \text{data parallelism beats the pipeline at layer } \ell
      &\iff \Delta_\ell \;<\; \max\!\Big(1,\; \frac{I^{*}}{I_\ell}\Big).
  \end{aligned}\;}
  \label{eq:cm-rule}
\end{equation}
Right of the ridge ($I_\ell > I^{*}$) the pipeline wins whenever data parallelism does more work,
and gsplit can lose only by the overlap gap. Left of the ridge, data parallelism wins as long as
the redundancy ratio $\Delta_\ell$ is below $I^{*}/I_\ell$ (or $1 + I^{*}/I_\ell$ against gsplit).
\assume{The comparison assumes both strategies produce the same distinct rows,
$u^{\mathrm{dp}}_\ell \approx u^{\mathrm{split}}_\ell$; the two samplings are not identical.}

\paraq{A data-parallel hybrid takes data parallelism for a prefix of layers.}
With split point $k$, layers $\ell < k$ (those nearest the seeds) run data-parallel and layers
$\ell \ge k$ run split. Because $\delta^{\mathrm{dp}}_\ell$ grows with $\ell$, the rule of
Eq.~\eqref{eq:cm-rule} tends to favour data parallelism for the layers nearest the seeds and gsplit
for the outer ones. The hybrid is realisable when the layers that data parallelism wins form a
prefix $\ell < k$, and it beats both pure strategies when that prefix is neither empty nor the
whole model.

\begin{figure}[h]
  \centering
  \fbox{\parbox{0.85\linewidth}{\centering\small
  \expdiag{classic roofline. $x$: FLOPs per interconnect byte $I$ (log); $y$: useful FLOP/s per
  GPU (log). Flat: dense roof $\hat{P}$ and gather roof $\hat{B}_{\mathrm{mem}}/g$ (one per model
  and $h$); sloped: $\hat{B}_{\mathrm{net}} I$ for PCIe-emu and NVLink. Ridge $I^*$ marked per
  interconnect. One colour per layer $\ell$: its pipeline point on the roof, its gsplit point
  directly below, and its data-parallel layer as a horizontal line at $C/\delta^{\mathrm{dp}}_\ell$.
  The winner of layer $\ell$ is the highest of the three.}\\[4pt]
  \placeholder}}
  \caption{Schematic classic roofline (placeholder).}
  \label{fig:cm-schematic}
\end{figure}

% ---------------------------------------------------------------------------
\subsection{Regions of the roofline}
\label{sec:cm-regions}

\paraq{The ridge divides the plane into two regions.}
\begin{center}
\footnotesize
\begin{tabular}{l l l l}
\toprule
Region & Condition & Bottleneck & Favours \\
\midrule
Communication-bound & $I < I^{*}$ & interconnect & DP if Eq.~\eqref{eq:cm-rule} holds, else split \\
Compute-bound & $I > I^{*}$ & dense arithmetic or gather & split \\
\bottomrule
\end{tabular}
\end{center}
Near the ridge both resources matter, and the pipeline's gain over gsplit is largest
(Eq.~\eqref{eq:cm-pipe-gain}). In the compute-bound region, whether the dense or the gather roof is
lower says which kernel a faster implementation would have to improve. It does not change the
strategy.

\paraq{Hardware moves the roofs; the workload moves the points.}
\begin{itemize}\itemsep1pt
  \item \textbf{Interconnect or GPU.} Changes $\hat{B}_{\mathrm{net}}$, $\hat{P}$, or
        $\hat{B}_{\mathrm{mem}}$, and with them the roofs and the ridge
        $I^{*}$ (Eq.~\eqref{eq:ridge}). The points stay where they are. A layer that is
        compute-bound on NVLink can become communication-bound on emulated PCIe.
  \item \textbf{Hidden size.} When the $h^2$ term of $F$ dominates, $I$ grows roughly linearly in
        $h$ and the point moves right; $g$ falls, so the gather roof rises.
  \item \textbf{Model.} More dense work per row (GAT) moves the point right of a SAGE layer of the
        same $h$, and GAT's gather roof differs from SAGE's.
  \item \textbf{Depth.} Outer layers have larger frontiers: $\delta^{\mathrm{dp}}_\ell$ grows and
        $m_\ell$ changes with the partition's cut. Different layers of one model can therefore
        fall on different sides of the rule~\eqref{eq:cm-rule}---the condition under which a
        per-layer hybrid beats both pure strategies.
\end{itemize}

% ---------------------------------------------------------------------------
\subsection{Assumptions}
\label{sec:cm-assumptions}

\paraq{All assumptions of the roofline in one place.}
Besides the roofline assumptions (R1)--(R5) of Section~\ref{sec:cm-purpose}, the model makes
these:
\begin{enumerate}\itemsep1pt
  \item[\assume{(M1)}] \assume{Uniform hidden width $h$; input-feature loading, the first layer's
        width, and sampling are outside the plot.}
  \item[\assume{(M2)}] \assume{Per-edge dense products charged per destination (over-counts
        GAT).}
  \item[\assume{(M3)}] \assume{The busiest GPU represents the layer.}
  \item[\assume{(M4)}] \assume{$u^{\mathrm{dp}}_\ell \approx u^{\mathrm{split}}_\ell$ when
        comparing strategies at one layer.}
  \item[\assume{(M5)}] \assume{The pipeline realises (R1) exactly.}
\end{enumerate}

% --- Kept for the author (not typeset): what changed relative to the earlier draft (roofline.tex, in git history at 8869a0ac~).
% Earlier draft                                  | This version                         | Simplification
% gsplit: compute + shuffle (sum); pipeline: max | all overlap (R1); gsplit = gap below | one roof formula; pipeline gain = gap below roof
% Axes in useful bytes; x = compute/useful byte  | x = FLOPs per interconnect byte      | classic axes; x independent of GPU
% Pull ratio r, per-layer ridge x*(r)            | r absorbed into I                    | one ridge I* per GPU and interconnect
% Dense/gather roofs as parallel sloped lines    | two flat ceilings, C = min(P, Bmem/g)| same information, standard picture
% Redundancy delta for DP only; busiest GPU      | one work inflation delta, all strats | one rule y = min(.)/delta
% Band parameter beta_band, four regions         | two regions split by the ridge       | no free parameter
% Model constant alpha (rows gathered per edge)  | renamed kappa (alpha = GAT attention)| no symbol clash

\paraq{Conclusion of this section.}
Each layer of each strategy is one point on a classic roofline. Its intensity $I$ and work
inflation $\delta$ come from sampling alone; the roofs come from three measured rates. Pipeline
points lie on the roof, gsplit points lie up to $2\times$ below them, and data-parallel layers are
flat lines lowered by redundancy. The boxed rule~\eqref{eq:cm-rule} names the winner of each
layer, and changing the GPU or the interconnect redraws the roofs over the same points.
Section~\ref{sec:predicted} draws this roofline for representative settings.
